% New lemmas addressing review items A, B, and the completion step.

\begin{lemma}[Jordan neighborhoods are local pairs]\label{lem:local-pair}
Let $A\subset\mathbb C$ be a nondegenerate full continuum compactly
contained in a Jordan disk $D_A$.  Then $D_A\setminus A$ is a
topological annulus; in particular $(D_A,A)$ is a local pair in the
sense of the localization construction
(Definition~\ref{def:localization}).
\end{lemma}

\begin{proof}
A connected component of $D_A\setminus A$ whose closure does not meet
$\partial D_A$ would be a bounded component of $\mathbb C\setminus A$,
contradicting fullness of $A$; hence $D_A\setminus A$ is connected.
In the sphere $\widehat{\mathbb C}$, the complement of $D_A\setminus A$
has exactly the two components $A$ and
$\widehat{\mathbb C}\setminus D_A$, so $D_A\setminus A$ is a doubly
connected plane domain, i.e.\ an annulus.  Its fundamental group is
infinite cyclic, generated by a core curve, which is the defining
property of a local pair.
\end{proof}

\begin{definition}[Shared local neighborhood and based subgroup
representatives]\label{def:shared-neighborhood}
Fix once and for all a Jordan disk $D_A\supset A$ with
$\overline{D_A}\subset U^{p}$ and
$\overline{D_A}\cap\bigl(F^{-1}B\setminus A\bigr)=\varnothing$; such a
disk exists because $F^{-1}B\setminus A$ is a compact set disjoint from
the component $A$ of $F^{-1}B$.  Since
$\Omega_s=T'\setminus A\subset T\setminus A=\Omega_t$ and
$D_A\setminus A\subset\Omega_s$, the same disk is a valid local
neighborhood on both sides.  With a common basepoint and connecting
path fixed in $D_A\setminus A$, define
\[
 H_s=\operatorname{im}\bigl[\pi_1(D_A\setminus A)\to\pi_1(\Omega_s)\bigr],
 \qquad
 H_t=\operatorname{im}\bigl[\pi_1(D_A\setminus A)\to\pi_1(\Omega_t)\bigr].
\]
These are the based representatives of the localization subgroups used
throughout; the collision certificate is representative-independent
(the winding homomorphism kills every conjugate of $H_s$).
\end{definition}

\begin{lemma}[Injectivity of the completion]\label{lem:completion-injective}
In the connected-complement localization construction
(Theorem~\ref{thm:kahn-localization}), the natural map from the
complement covering surface $\widehat\Omega$ into the completed
localized surface $S_A$ is injective.
\end{lemma}

\begin{proof}
The completed surface is the quotient of the disjoint union
$U\sqcup\widehat\Omega$ in which each point $u\in U\setminus A$ is
identified with its image $\ell(u)\in\widehat\Omega$ under the
canonical lift $\ell$ of $U\setminus A$.  Every identification therefore
pairs one point of $U$ with one point of $\widehat\Omega$; no two
points of $\widehat\Omega$ are ever identified with one another, since
$\ell$ is single-valued.  Hence the composition
$\widehat\Omega\hookrightarrow U\sqcup\widehat\Omega\to S_A$ is
injective.
\end{proof}

\begin{remark}\label{rem:completion-consequence}
Lemma~\ref{lem:completion-injective}, applied on both sides, shows that
the filled completion identifies no pair of distinct complement-cover
points.  Consequently every point collision of the completed localized
inclusion $\widehat\imath$ arises from a collision of the
complement-cover lift $\widehat a$, and conversely; this is the bridge
invoked in Theorem~\ref{thm:forced-collision}, now proved rather than
delegated to the verification checklist of
Section~\ref{sec:verification}.
\end{remark}
